\section{拟合外推及讨论}
\subsection{条件响应代理与分组验证}
预测目标为式\eqref{eq:paper-dq}的$\Delta Q$和式\eqref{eq:paper-margin}的$\Delta M$。对每个位值与目标分别建立模型，共四个回归器。特征为
\begin{equation}
\mathbf x=\left(Q_{\rm dep},\;Q_{\rm dep}/L,\;Q_{\rm dep}\sin\theta\right),
\label{eq:paper-features}
\end{equation}
依次采用fC、$\mathrm{fC}/\mu\mathrm m$、fC为单位；模型使用训练折均值和总体标准差标准化。$Q_{\rm dep}/L$引入尺寸相关尺度，角度项提供简单方向依赖，但这些组合是低维代理特征，不构成对微观收集机制的唯一解释。

岭回归最小化
\begin{equation}
\sum_{i\in\mathcal T}\left[y_i-\beta_0-\boldsymbol\beta^{\mathsf T}\widetilde{\mathbf x}_i\right]^2
+\lambda\|\boldsymbol\beta\|_2^2,
\label{eq:paper-ridge}
\end{equation}
其中截距不受惩罚。候选$\lambda$为$0$、$0.01$、$0.1$、$1$、$10$。训练数据按$(L,\text{粒子},\theta)$分成八个物理工况，逐组留出；每组四次重复及两种位值保持同组，避免同事件跨训练与验证。选择指标是八折RMSE的算术平均，四个模型均选得$\lambda=0.1$。标准化只使用该折训练数据，最终拟合才使用全部训练响应。比较基准是对应训练折目标均值。

\subsection{训练交叉验证结果及其解释}
表\ref{tab:paper-training-cv}和图\ref{fig:paper-training-cv}给出训练阶段结果。存储0/1的$\Delta Q$平均折RMSE为$0.102888/0.112903\fC$，训练均值基准为$2.279967/2.515710\fC$；$\Delta M$为$0.955160/1.012491\mV$，均值基准为$21.458870/23.383861\mV$。在当前工况中，引入沉积、尺寸和角度特征比常数预测保留了更多响应变化信息。

\begin{table}[htbp]
\centering\small
\caption{八组物理工况留一交叉验证的平均折RMSE；属于训练阶段模型选择，不是独立测试误差}
\label{tab:paper-training-cv}
\begin{tabular}{clrr}
\toprule
存储位 & 目标（单位） & 岭回归 RMSE & 均值基准 RMSE \\
\midrule
0 & $\Delta Q$ (fC) & 0.102888 & 2.279967 \\
0 & $\Delta M$ (mV) & 0.955160 & 21.458870 \\
1 & $\Delta Q$ (fC) & 0.112903 & 2.515710 \\
1 & $\Delta M$ (mV) & 1.012491 & 23.383861 \\
\bottomrule
\end{tabular}

\end{table}
\begin{figure}[htbp]
\centering
\includegraphics[width=0.9\textwidth]{r50-training-cv.pdf}
\caption{训练交叉验证中岭回归与训练均值基准的比较。不同目标分别保留fC和mV单位；误差来自八个工况折的RMSE平均，不能当作全部留出点合并后的RMSE或独立外推精度。}
\label{fig:paper-training-cv}
\end{figure}

上述误差是用于选择正则化参数的同一交叉验证过程的结果，不能视为无偏的最终泛化估计。特别是八折留出一个物理工况，并不等同于完整留出一个从未参与训练的栅长。误差低于均值基准支持该参数化表示的训练内可行性；它尚未回答$0.25\um$插值和$0.45\um$外推是否可靠，也未验证预测分类能否稳定越过$5\mV$边界。

\subsection{独立预测的冻结条件与待验收范围}
模型、标准化参数以及64个测试位响应的预测已在独立测试响应提交前冻结。最终评分必须复用这些预测，分别给出插值留出与尺寸外推的MAE、RMSE、最大绝对误差及分层残差，不按测试表现重新调参，也不筛去不利工况。设计未规定必须优于均值基准才纳入结果；真实完成且通过对应审计的测试应按预定规则报告，包括预测失败或优势不明显的情况。

本代理的输入包含真实Geant4沉积，目标的定义和绝对裕量还依赖已知无粒子基线，因此它是给定源项和基线条件下的TCAD响应代理。它没有从任意粒子描述直接预测器件结果，也没有替代所有Geant4与无粒子TCAD计算。本文快照尚未完成独立测试全量验收，不给出其最终误差、预测对真值散点图或外推通过结论。

\subsection{模型适用边界与后续接入}
尺寸变化可能同时改变输运、端口耦合和电荷共享动态，本模型只改变栅长并保持多项参数固定，无法分离真实工艺节点缩放中的全部因素。当前质子结果限于所选电磁配置下的电离响应，不能据此讨论核反应次级粒子或核反冲主导的单粒子效应。硅/真空几何、等效直径迹和二维归一化也未覆盖完整材料栈及真实三维敏感体积。源积分守恒和数值一致性是必要检查，但不能替代实验或可信工艺模型的校准。

已规划的$4/6\,\mathrm{nm}$氧化层比较采用固定硅入口事件，只能作为TCAD参数敏感性分析；它不构成跨完整工艺模型的粒子输运重算，也不属于当前栅长预测的氧化层外推验证。其余扩展、边界数值检查及R50原生图形取证尚待各自完成归档，不在本文画占位结果。后续结果应保持R50版本、事件身份和匹配条件一致，并保留失败及确定性重跑的对应关系，避免把不同模型或重复求解混成额外物理证据。
